% !TEX root = ../main.tex

%----------------------------------------------------------------------------------------
% CHAPTER 1
%----------------------------------------------------------------------------------------

% LNW3-COUNT-START
\vspace*{-3em}
\chapter{Einleitung}

Gemeinnützige Organisationen finanzieren ihre Arbeit überwiegend aus Spenden und leisten damit einen wesentlichen Beitrag zur Bewältigung sozialer und ökologischer Herausforderungen. Ihre Wirkung hängt unmittelbar davon ab, ob potenzielle Spenderinnen und Spender den Organisationen vertrauen. Aktuelle Studien zeigen jedoch, dass intransparente Mittelverwendung als einer der zentralen Gründe für sinkendes Vertrauen und damit für eine rückläufige Spendenbereitschaft gilt~\cite{SUNG2023103812}.

Als technischer Lösungsansatz wird zunehmend die Blockchain-Technologie diskutiert, weil sie Transaktionen öffentlich nachvollziehbar macht und vermittelnde Institutionen prinzipiell ersetzen kann~\cite{Teng2021}. Sie verspricht damit, das Vertrauensdefizit auf Infrastruktur-Ebene strukturell zu adressieren.

Der vorliegende Bericht befasst sich vor diesem Hintergrund mit der Plattform \emph{The Giving Block}, einem etablierten Anbieter für Kryptowährungsspenden an gemeinnützige Organisationen~\cite{givingblock}. Ziel des Berichts ist darzustellen, wie die Plattform Blockchain-Mechanismen einsetzt, um das beschriebene Vertrauensproblem im Spendenwesen zu adressieren, und welcher gesellschaftliche Nutzen daraus entstehen kann. Kapitel~\ref{chapter2} ordnet hierzu den Stand der Technik blockchainbasierter Spendensysteme ein, Kapitel~\ref{chapter3} beschreibt den Lösungsansatz von \emph{The Giving Block} und Kapitel~\ref{chapter4} bewertet die Ergebnisse mit Blick auf den praktischen und gesellschaftlichen Nutzen.


\chapter{Ausgangslage}
\label{chapter2}

\section{Problemstellung}
Klassische Vertrauensinstrumente im Spendensektor wie Jahresberichte, Spendensiegel oder externe Audits wirken überwiegend retrospektiv und setzen vermittelnde Institutionen voraus. Empirische Untersuchungen zeigen, dass dieses indirekte Modell den heutigen Transparenzerwartungen nicht mehr genügt: Spendende fordern Nachvollziehbarkeit auf Transaktionsebene, was klassische Berichtsformate strukturell nicht leisten können. Daraus ergibt sich die Leitfrage, wie eine technologisch gestützte Plattform diese strukturelle Lücke konkret schliessen kann.

\section{Stand der Technik}
Als technische Antwort auf das beschriebene Transparenzdefizit wird die Blockchain-Technologie diskutiert. Eingeführt durch Nakamoto~\cite{nakamoto2008}, basiert sie auf einer Kette digital signierter Transaktionen. Abbildung~\ref{fig:blockchain-tx} zeigt diese Verkettung schematisch: Jeder Besitzer überträgt eine Einheit, indem er den Hash der vorherigen Transaktion und den öffentlichen Schlüssel des nächsten Empfängers mit seinem privaten Schlüssel signiert. So entsteht eine kryptografisch verkettete Eigentumshistorie, die von jedem Teilnehmer ohne zentrale Instanz verifiziert werden kann.

\begin{figure}[h]
    \centering
    \includegraphics[width=0.8\textwidth]{transaction.png}
    \caption{Transaktionskette als Folge digitaler Signaturen~\cite{nakamoto2008}.}
    \label{fig:blockchain-tx}
\end{figure}

Auf dieser Grundlage weisen Sung et al. experimentell nach, dass blockchainbasierte Spendensysteme über Transparenz, Unveränderbarkeit und Effizienz die wahrgenommene Vertrauenswürdigkeit von NPOs signifikant erhöhen~\cite{SUNG2023103812}. Eine US-Studie mit 590 Befragten bestätigt entsprechende Effekte auf Einstellungen und subjektive Normen der Spendenden, moderiert allerdings durch Faktoren wie Religion und persönliches Risikoempfinden~\cite{BAUDIER2023122524}. Trotz dieses Potenzials bleibt die Adoption begrenzt: Die On-Chain-Transparenz endet, sobald Mittel in Fiat-Währungen umgewandelt werden, die Kursvolatilität schafft planerische Risiken, und die Bedienung bildet für nicht-technische Spendende eine Akzeptanzhürde. Abbildung~\ref{fig:vergleich-transparenz} verdeutlicht, dass Blockchain vor allem Granularität und Zeitbezug der Transparenz verändert.

\begin{figure}[h]
    \centering
    \begin{tabularx}{\textwidth}{@{}lXX@{}}
        \toprule
        \textbf{Dimension} & \textbf{Klassisches Modell} & \textbf{Blockchainbasiertes Modell} \\
        \midrule
        Verifikationsinstanz & Externer Auditor / Spendensiegel & Konsensverfahren des Netzwerks \\
        Granularität         & Aggregierte Jahresberichte       & Einzeltransaktion (on-chain) \\
        Zeitbezug            & Retrospektiv                     & Echtzeit \\
        Vertrauensbasis      & Institutionell vermittelt        & Kryptografisch garantiert \\
        Reichweite           & National (NGO-Sitzland)          & Geografisch entgrenzt \\
        \bottomrule
    \end{tabularx}
    \caption{Vergleich klassischer und blockchainbasierter Spendentransparenz (eigene Darstellung).}
    \label{fig:vergleich-transparenz}
\end{figure}


\chapter{Lösungsansatz}
\label{chapter3}

\emph{The Giving Block} adressiert die in Kapitel~\ref{chapter2} beschriebene Umsetzungslücke, indem die Plattform Kryptowährungsspenden für gemeinnützige Organisationen technisch und organisatorisch standardisiert~\cite{givingblock}. Sie agiert als Zahlungsinfrastruktur auf öffentlichen Distributed Ledgers und bündelt Wallet-Verwaltung, optionale Fiat-Konvertierung, Spendenbelege, Reporting und Compliance-Werkzeuge. Dadurch müssen NGOs keine eigene Blockchain-Infrastruktur betreiben und können Krypto-Spenden dennoch in Buchhaltungs- und Steuerprozesse einordnen.

Funktional initiiert die spendende Person eine Krypto-Transaktion an die Wallet-Adresse der Organisation; das Netzwerk verifiziert sie über asymmetrische Kryptografie und ein Konsensverfahren und protokolliert sie unveränderbar im Distributed Ledger~\cite{nakamoto2008, MeierStormer2018}. Über öffentliche Block-Explorer können alle Beteiligten den Transfer jederzeit nachprüfen, womit Nachvollziehbarkeit nicht mehr institutionell vermittelt, sondern technisch garantiert wird~\cite{Teng2021}. Abbildung~\ref{fig:spendenfluss} fasst diesen Ablauf zusammen.

\begin{figure}[h]
    \centering
    \begin{tikzpicture}[
        node distance=1.2cm and 2.6cm,
        block/.style={draw, rounded corners=2pt, minimum width=2.2cm, minimum height=0.95cm, align=center, font=\small, fill=zhawblue!10},
        ledger/.style={draw, rounded corners=2pt, minimum width=2.4cm, minimum height=0.95cm, align=center, font=\small, fill=zhawblue!22},
        explorer/.style={draw, rounded corners=2pt, minimum width=4.4cm, minimum height=0.85cm, align=center, font=\small, fill=black!8},
        arrow/.style={-{Stealth[length=2.2mm]}, thick},
        lbl/.style={font=\footnotesize, fill=white, inner sep=1pt}
    ]
        \node[block] (donor) {Spender:in};
        \node[ledger, right=of donor] (ledger) {Distributed\\Ledger};
        \node[block, right=of ledger] (ngo) {NGO};
        \node[explorer, below=of ledger] (explorer) {Öffentlicher Block-Explorer\\(jederzeit verifizierbar)};

        \draw[arrow] (donor) -- node[above, lbl]{1.~Krypto-TX} (ledger);
        \draw[arrow] (ledger) -- node[above, lbl]{3.~Auszahlung} (ngo);
        \draw[arrow] (ledger) -- node[right, lbl]{2.~Konsens} (explorer);
    \end{tikzpicture}
    \caption[Spendenfluss über The Giving Block]{Vereinfachter Ablauf einer Krypto-Spende: Initiierung durch Spender:in, Konsens-Verifikation und unveränderbare Protokollierung im Distributed Ledger, öffentliche Nachvollziehbarkeit über Block-Explorer und abschliessende Mittelweitergabe an die NGO mit optionaler Fiat-Konvertierung (eigene Darstellung).}
    \label{fig:spendenfluss}
\end{figure}

In der Schweiz gelten Kryptowährungen steuerlich als bewegliches Vermögen, weshalb Krypto-Zuwendungen analog zu Sachspenden behandelt werden und unter den üblichen Voraussetzungen abzugsfähig bleiben~\cite{ESTV2022}. Damit übersetzt \emph{The Giving Block} die in Kapitel~\ref{chapter2} dokumentierte Vertrauenswirkung blockchainbasierter Spendensysteme in eine praktikable Spendenpraxis: Transparenz entsteht on-chain, während Quittierung, Umrechnung und Compliance über die Plattform administrierbar bleiben. Die Reichweite hängt jedoch weiterhin von Bedienbarkeit, Akzeptanz und der Transparenz nach einer Fiat-Konvertierung ab.



\chapter{Fazit}
\label{chapter4}

\emph{The Giving Block} beantwortet die Leitfrage des Berichts, indem die Plattform institutionell vermitteltes Vertrauen durch kryptografisch garantierte Nachvollziehbarkeit ergänzt und in einen steuerlich anschlussfähigen Spendenkreislauf einbettet. Gesellschaftlich kann dies Vertrauen in NGOs stärken, internationale Spenden erleichtern und digital affine Zielgruppen erreichen. Grenzen bleiben bei der Fiat-Konvertierung und der Spenderakzeptanz. Stablecoins und Smart-Contract-basierte Verwendungsnachweise könnten diese Wirkung künftig verstärken.

% LNW3-COUNT-END
\newpage
\vspace*{-3em}
\chapter{Verzeichnisse}

\section*{Literaturverzeichnis}
\addcontentsline{toc}{section}{Literaturverzeichnis}
\printbibliography[heading=none]

\section*{Abbildungsverzeichnis}
\addcontentsline{toc}{section}{Abbildungsverzeichnis}
\makeatletter\@starttoc{lof}\makeatother



\chapter{Deklaration KI}
\label{chapter5}

Im Verlauf dieser Arbeit kamen generative KI-Systeme in mehreren Prozessphasen zum Einsatz. Im Einzelnen wurde \emph{Claude} (Anthropic) für die Strukturierung und Formulierungsverbesserung einzelner Textpassagen, für das Aufsetzen und Anpassen der \LaTeX- und BibTeX-Vorlagen sowie für das Debugging von Layout- und Kompilationsproblemen verwendet. Die in Claude integrierte Web-Suchfunktion wurde ergänzend für die Recherche zusätzlicher deutschsprachiger Fachquellen (insbesondere Quellen~\cite{MeierStormer2018} und~\cite{ESTV2022}) eingesetzt; die bibliografischen Angaben dieser Quellen wurden anschliessend manuell verifiziert.

Die inhaltliche Auswahl der Themen und Quellen, die fachliche Prüfung sämtlicher Aussagen sowie die abschliessende Redaktion des Berichts erfolgten durch die Autoren. Verantwortung für Relevanz, Wahrheitsgehalt und korrekte Quellenangabe liegt entsprechend Ziff.~4 der ZHAW-Richtlinie zur Deklarationspflicht von generativer KI bei Arbeiten an der School of Engineering vollständig bei den Autoren.
